\section{Related Work}

\subsection{Remote Photoplethysmography (rPPG) Extraction}
Classical rPPG algorithms focus on recovering the underlying blood volume pulse (BVP) from noisy RGB camera signals. Techniques such as Independent Component Analysis (ICA) and Chrominance-based methods (CHROM) paved the way for robust spatial projection by isolating the pulsatile subspace from lighting variations. Notably, Wang et al. introduced the Plane-Orthogonal-to-Skin (POS) algorithm \cite{wang2016pos}, which elegantly cancels specular reflections by projecting facial color signals onto a plane orthogonal to the localized skin-tone vector. Recent deep learning approaches have attempted to extract BVP waveforms directly using 3D Convolutional Neural Networks (CNNs) and attention-based temporal shift networks (TS-CAN) \cite{liu2020tscan}. However, pure end-to-end video models often struggle to maintain physiological constraints, frequently hallucinating pulse waveforms that lack true hemodynamic meaning, and are highly susceptible to overfitting on small datasets.

\subsection{Deep Learning for Cuffless Blood Pressure Estimation}
Extracting blood pressure from photoplethysmogram signals is an active area of research. Traditionally, this domain relied on handcrafted morphological features such as Pulse Transit Time (PTT) and Pulse Wave Velocity (PWV). While theoretically sound, PTT-based methods strictly require dual-sensor modalities (e.g., synchronous ECG and PPG measured at distal sites), rendering them impractical for ubiquitous single-camera facial monitoring. 

To overcome this, deep learning has shifted the paradigm toward single-channel feature extraction. Recurrent architectures like Long Short-Term Memory (LSTM) networks and Gated Recurrent Units (GRUs) are commonly employed to model the temporal dependencies and implicit morphological landmarks within a single pulsatile waveform \cite{slapnicar2019blood}.

Despite these advances, existing rPPG-to-BP frameworks face significant limitations. When confronted with heavily damped inputs—inherent to facial microvasculature—monolithic deep learning architectures predicting both Systolic BP (SBP) and Diastolic BP (DBP) concurrently often suffer from the "regression-to-the-mean" phenomenon \cite{shradha2023regression}. Because SBP and DBP vary independently, updating a shared feature representation with conflicting gradients forces the model to collapse toward a narrow, normotensive baseline band (e.g., predicting 120/80 mmHg for all subjects) to minimize aggregate Mean Squared Error. This leads to catastrophic failure when predicting hypertensive crises or hypotensive shock. Our work directly addresses these challenges by introducing Savitzky-Golay morphological preservation and a mathematically decoupled regression architecture (MODEL-06-SepHead) to isolate gradient flows.
